\documentclass{article}
\usepackage[paperwidth=24.25cm, paperheight=10.29cm,left=0pt,top=0pt,bottom=0pt,right=0pt]{geometry}
\usepackage[french]{babel}
\usepackage[autolanguage]{numprint}
\input{./figspreamble}
\usepackage{tikz}
\begin{document}\noindent
\colorlet{goodcolor}{green!60!black}%
\colorlet{badcolor}{orange!90!black}%
\colorlet{failcolor}{red}%
\tikzset{%
assurance/.style={text width=5cm, align=center,anchor=north,thick,rounded corners=5pt,minimum height=50pt},
faute/.style={draw=#1,fill=#1!30},
good/.style={goodcolor},
bad/.style={badcolor},
fail/.style={failcolor}
}%
\begin{tikzpicture}[
  row 1/.style={draw,text width=5cm,align=center,anchor=north},
  row 2/.style={assurance}, 
  every path/.style={-stealth,thick},
  level/.style={sibling distance=5cm},
  column 2/.style={opacity=0.3},
  column 3/.style={opacity=0.3}
]
\node[matrix] {
\node {RESPONSABILITÉ PÉNALE}; &[5mm] \node {RESPONSABILITÉ CIVILE\\(OBLIGATOIRE)}; &[2mm] \node {ASSURANCE INDIVIDUELLE\\(OPTIONNELLE)}; \\
\node {Non-respect de la loi (dont Code du Sport). Baptême au-delà de \numprint{6}~mètres qui tourne mal, \dots}
  child { node[faute=failcolor] {Infraction pénale}
    child[fail] {
      node[text width={}] (juge) {Juge, tribunal,\\amende, prison, \dots}
      child[fail] {
        node[fail,text width={}] (noass) {Zéro assurance}
      }
    }
  };
&
\node at (0,0) {J'ai cassé l'appareil photo d'un collègue plongeur, \dots}
  child { node[faute=badcolor] {Dommage corporel, moral ou matériel causé à un tiers}
    [every node/.append style={text width={}}]
    child[bad] {
      node {Préjudice volontaire\\ou hors contrat} child {node {Pour ma pomme}}
    }
    child[good] {
      node {Préjudice involontaire\\garanti par contrat} child { node {Remboursé en partie\\ou totalité par l'assurance\\en responsabilité civile (RC)}}
    }
  };
&
\node at (0,0) {J'ai eu un accident quelconque qui nécessite un rapatriement, des soins, \dots}
  child { node[faute=badcolor] {Dommage corporel sans tiers responsable}
    [every node/.append style={text width={}}]
    child[bad] {
      node {Je n'ai pas\\d'assurance individuelle} child {node {Pour ma pomme}}
    }
    child[good] {
      node {J'ai une assurance\\individuelle} child { node {Remboursé en partie\\ou totalité par l'assurance\\individuelle}}
    }
  };
\\
};
\end{tikzpicture}
\end{document}
